% =========================================================================
% secciones_zh/sec07_conclusion_future_work.tex — 第VII节 + 术语表 + 附录A
% =========================================================================
\section{结论与未来工作}
\label{sec:conclusion}

本文贡献了SENI-142：一个开放的、实测的、确定性可复现的完整国家输电系统基准——142节点、约112座变电站、135条线路、58台变压器、第368-22号法律下的十个领土微电网——其精确网络分块凝聚以2.65~ms与473~MiB（作者8线程研究工作站上的实验性参考测量）求解每个小时级国家级交流状态，相量残差 $\delta = 2.788 \times 10^{-14}$~pu，以十年连续运行遥测验证，并由四相数字孪生生命周期封存。基准、验证工具与韧性记录经IEEE DataPort与GitHub开放发布。

四条研究路线由此直接展开。第一，\textbf{亚秒级物理信息神经代理}：87,360状态水平域是经典扫描之学习加速器的训练场，精确求解器保留为证伪仲裁者\cite{misyris2020}。第二，\textbf{构网型变流器EMT联合仿真}：将构网（grid-forming）资源的变流器控制级动态\cite{rosso2021}、\cite{ieee2800}嵌入分块包络，以研究频率的形成而不仅是频率的支撑。第三，\textbf{暂态调度下电解槽膜老化}：电转气区块的跨时段运行正是PEM电堆老化研究的天然负荷曲线\cite{irena2020}、\cite{iea2023}。第四，\textbf{硅上分布式小时结算}：结算级小时水平域、签名清单与默克尔式溯源链\cite{merkle1987}、\cite{fips180}勾勒出一种结算架构，使每个计价小时都与密码学证明的物理状态对账。这座无法输入惯量的岛屿，在自身的实测镜像中产出了可输出的东西：一份面向能源自主工程的公共仪器。

% ---- 术语表（英文版内容的表格形式）----
\begin{table*}[!t]
\caption{术语表}
\label{tab:nomenclature}
\centering
\footnotesize
\begin{tabular}{@{}p{0.24\textwidth} p{0.64\textwidth}@{}}
\toprule
$\mathcal{G} = (\mathcal{V}, \mathcal{E})$ & 输电网络的连通无向图 \\
$\mathcal{V}$，$|\mathcal{V}| = 142$ & 输电节点集 \\
$\mathcal{E}_{\mathrm{line}}$，$\mathcal{E}_{\mathrm{trafo}}$ & 线路支路（135），变压器支路（58） \\
$\mathbf{Y}_{\mathrm{bus}}$ & 节点导纳矩阵，$\mathbb{C}^{142 \times 142}$ \\
$S_i = P_i + j Q_i$ & 节点 $i$ 的复功率注入 \\
$Z_{\mathrm{cc}}$ & 被化简区域的短路阻抗 \\
$\mathcal{G}_k$，$\Omega_k$ & 领土微电网 $k$ 及其宏域 \\
$\mathcal{V}_B$，$\Gamma$ & 边界节点集与凝聚边界 \\
$\mathbf{A}, \mathbf{B}, \mathbf{C}, \mathbf{D}$ & 重排后系统矩阵的内部/边界分块 \\
$\mathbf{S} = \mathbf{D} - \mathbf{C}\mathbf{A}^{-1}\mathbf{B}$ & 内部块的舒尔补 \\
$\mathbf{x}_I$，$\mathbf{x}_B$，$\mathbf{u}_I$，$\mathbf{u}_B$ & 内部/边界状态与右端项 \\
$N$，$N_k$，$N_B$ & 系统、子域与边界维度 \\
$\kappa_2(\cdot)$ & 谱条件数 \\
$\delta$ & 相量边界耦合残差（pu） \\
$P_i(\mathbf{V}_t, \boldsymbol{\theta}_t)$ & 由电网状态得到的节点有功注入 \\
$P^{\mathrm{EL}}_{h,t}$，$P^{\mathrm{FC}}_{h,t}$ & 机组 $h$ 的电解槽/燃料电池电功率 \\
$\eta_{\mathrm{EL}}$，$\eta_{\mathrm{FC}}$ & 电解/再转换效率 \\
$E_{h,t}$，$E_h^{\min}$，$E_h^{\max}$ & 氢缓冲库存及其上下界 \\
$\Delta t$ & 时间步长（1 h） \\
$C_p(\lambda, \beta)$ & 风机功率系数 \\
$\sigma(\cdot)$，$\alpha$ & S形激活及其自适应斜率 \\
$Q_g^{\lim}$，$Q_g^{\mathrm{eff}}$ & 发电机无功限值与正则化能力 \\
PHES & 抽水蓄能（764 MW可逆机组群） \\
ADM1 & 一号厌氧消化模型 \\
PEM & 质子交换膜（电解） \\
\bottomrule
\end{tabular}
\end{table*}

% ---- 附录A（技术参数表，精修指令）----
\appendix
\section{SENI-142岛屿基准的拓扑与电气汇总参数}
\label{app:parameters}

表~\ref{tab:parameters}汇总已发布基准在节点模型中的结构与电气参数，使第~\ref{sec:topology}--\ref{sec:benchmarking}节的拓扑算术可对照已发布数据集一目了然地核验。

\begin{table*}[!t]
\caption{SENI-142岛屿基准的拓扑与电气汇总参数。上半部分为本工作从CAD单线图恢复的领土实例化，构成本文自身的实测贡献；下半部分为所有已报告仿真实际运行的电气计算模型。}
\label{tab:parameters}
\centering
\footnotesize
\begin{tabular}{@{}l r@{}}
\toprule
\textbf{量} & \textbf{数值} \\
\midrule
\multicolumn{2}{@{}l}{\textit{领土实例化（本工作，实测）}} \\
主干输电变电站 & 112 \\
主干输电母线，345~kV & 6 \\
主干输电母线，230~kV & 3 \\
主干输电母线，138~kV & 106 \\
由69~kV凝聚的边界端口 & 27 \\
\textbf{节点计算母线总数} & \textbf{142} \\
\midrule
\multicolumn{2}{@{}l}{\textit{SENI-142电气计算模型（推导，凝聚）}} \\
输电线路 & 135 \\
电力变压器 & 58 \\
装机发电容量（铭牌） & 5,420 MW \\
最大小时需求（SCADA记录） & 3,450 MW \\
\bottomrule
\end{tabular}
\end{table*}

节点普查遵循第~\ref{sub:substations}节的节点规则：每座围栏场站每个电压等级一个节点，变压器组为支路。表~\ref{tab:parameters}上半部分为本工作从CAD单线图恢复的领土实例化，即本文自身的实测贡献：112座主干输电变电站之上的142个节点计算母线，其中115条为345、230与138~kV主干母线，27个为69~kV层凝聚而成的边界端口。下半部分为求解器实际运行的电气计算模型——该142母线电网之上的135条输电线路与58台电力变压器——其中69~kV次输电层由第~\ref{sub:reduction}节保持注入的节点等效表示。两个区块分开报告且是刻意的：上为测量，下为建模决策。第~\ref{sub:dataprotocol}节与第~\ref{sub:clause}节的两条逐字协议条款仍是数据层与领土实例化的契约核心，在英文规范版中逐字保留。
